\clearpage\thispagestyle{empty}\enlargethispage{2\baselineskip}\begingroup\small\raggedright
\noindent{\Large\bfseries Quellen und Terminologie\par}\medskip
Historischer Text: Carl Friedrich Gauss, \emph{Werke}, dritter Band, Göttingen, 1866. Diese englischen und deutschen Texte sind aus der originalsprachlichen Ausgabe KI-übersetzt. Die folgenden menschlichen Übersetzungen und Schriften werden für die angegebene Terminologie und den Sprachgebrauch herangezogen; ihr fortlaufender Wortlaut wird nicht übernommen. Die ursprünglichen mathematischen Formulierungen von Gauss bleiben maßgeblich.
\par\medskip\textbf{Paul Taylor; Bernard Leak.} Taylors englische Übersetzung des zweiten Beweises von Gauss für den Fundamentalsatz der Algebra, mit Korrekturen von Leak, 1983. Die entsprechende Übersetzung des Beweises von 1815 dient der englischen Terminologie zu Polynomen, Teilern, Graden, symmetrischen Funktionen und Beweisen. Verwandte Terminologie wurde auch für die Dissertation von 1799, die Vorzeichenarbeit von 1828, die Einleitung und die Artikel 1 und 10 von 1849, die Anzeigen von 1815 und 1849 sowie die Anzeigen zu Coytier und Fourier herangezogen; zu diesen ist dies keine Parallelübersetzung. Fortlaufender Text aus dieser menschlichen Übersetzung wird nicht übernommen. Ihr späterer Gebrauch komplexer Größen ersetzt hier nicht die Definition imaginärer Größen von Gauss. Die herangezogene Online-Ausgabe nennt Dezember 1983 als Übersetzungsdatum, August 2003 für die LaTeX-Fassung und den 8. Juni 2007 als letztes Änderungsdatum der zugrunde liegenden Quelle. Ihre Danksagung nennt Mark Wainwright für die Rückgabe von Taylors Notizen und Nikita Danilov für die Korrektur von Tippfehlern.\par
\url{https://www.paultaylor.eu/misc/gauss.html}
\par\medskip\textbf{H. Maser.} Deutscher Übersetzer und Herausgeber, \emph{Carl Friedrich Gauss’ Untersuchungen über höhere Arithmetik}, Berlin: Julius Springer, 1889. Verwandte deutsche Terminologie für algebraische Funktionen, Wurzeln, Teilbarkeit, Quotienten, Grade, Logarithmen, Tafeln und den Sprachgebrauch zu Beweisen und Existenz; keine Parallelübersetzung dieser Arbeiten.\par
\url{https://archive.org/details/carlfriedrichga00gausgoog}
\par\medskip\textbf{James Caddall Morehead und Adam Miller Hiltebeitel.} Übersetzer von Gauss’ \emph{General Investigations of Curved Surfaces of 1827 and 1825}, Princeton University Library, 1902; Project Gutenberg 36856, TeX-Ausgabe. Verwandte englische Terminologie zu Geometrie, Trigonometrie, Reihen, Koordinaten, Stetigkeit, Integration, vorzeichenbehafteten Flächen und unendlich kleinen Größen; keine Parallelübersetzung dieser algebraischen Beweise.\par
\url{https://www.gutenberg.org/ebooks/36856}
\par\medskip\textbf{Ulrich Görtz.} \emph{Algebra}, Wintersemester 2021/22, Abschnitt 3.1.1, Beispiele 3.1–3.2. Verwandter zeitgenössischer deutscher Sprachgebrauch zu Polynomen, Unbestimmten, Summen, Zählern und Nennern sowie zur mathematischen Argumentation. Diese Quelle dient der Terminologie; sie ist keine Parallelübersetzung dieser Beweise. Fortlaufender Text und moderne ringtheoretische Argumente werden nicht übernommen.\par
\url{https://math.ug/a-ws2122/sect0006.html}
\par\medskip\textbf{KI-Urheberschaft und Quellen.} Mit OpenAI GPT-6.1 Sol KI-übersetzt; KI-erzeugter Satz. Gauss ist der Autor des historischen mathematischen Textes. Die menschlichen Quellen dienen der genannten Terminologie und dem Sprachgebrauch.
\par\medskip Bibliografische Metadaten: \texttt{CITATIONS.bib}.
\par\endgroup\clearpage
